\setcounterpageref{ex}{0}
\TaoPhieuKiemTra{15}

\setcounter{ex}{0}
\setcounter{bt}{0}

\textbf{\large I. PHẦN TRẮC NGHIỆM (3,0 điểm)}

{\itshape Khoanh tròn vào chữ cái đặt trước câu trả lời đúng. Mỗi câu đúng được 0,5 điểm.}

\begin{ex} % Câu 1
Số chẵn liền trước số $50$ là
\choice
{\True $48$}
{$49$}
{$52$}
{$40$}
\loigiai{
Số chẵn liền trước $50$ là $48$.\\
\textbf{Đáp án đúng: A}
}
\end{ex}

\begin{ex} % Câu 2
Nếu số hạng thứ nhất là $32$, số hạng thứ hai là $15$ thì tổng là
\choice
{$46$}
{\True $47$}
{$48$}
{$17$}
\loigiai{
Tổng: $32 + 15 = 47$.\\
\textbf{Đáp án đúng: B}
}
\end{ex}

\begin{ex} % Câu 3
Tính: $9\text{ kg} + 5\text{ kg} = \dots$
\choice
{$13\text{ kg}$}
{$15\text{ kg}$}
{\True $14\text{ kg}$}
{$4\text{ kg}$}
\loigiai{
$9 + 5 = 14\text{ kg}$.\\
\textbf{Đáp án đúng: C}
}
\end{ex}

\begin{ex} % Câu 4
Nối $3$ điểm không thẳng hàng bằng thước kẻ ta được hình gì?
\choice
{Hình tứ giác}
{Đường thẳng}
{Đường cong}
{\True Hình tam giác}
\loigiai{
Nối 3 điểm không thẳng hàng ta được hình tam giác.\\
\textbf{Đáp án đúng: D}
}
\end{ex}

\begin{ex} % Câu 5
So sánh: $9 + 8 \dots 10 + 6$.
\choice
{\True $>$}
{$<$}
{$=$}
{Không so sánh được}
\loigiai{
$9 + 8 = 17$; $10 + 6 = 16$. Vì $17 > 16$ nên điền dấu $>$.\\
\textbf{Đáp án đúng: A}
}
\end{ex}

\begin{ex} % Câu 6
Tính: $7 + 9 - 8 = \dots$
\choice
{$7$}
{\True $8$}
{$9$}
{$10$}
\loigiai{
$7 + 9 = 16$; $16 - 8 = 8$.\\
\textbf{Đáp án đúng: B}
}
\end{ex}

\vspace{4mm}
\textbf{\large II. PHẦN TỰ LUẬN (7,0 điểm)}

{\itshape Thí sinh trình bày chi tiết bài làm vào giấy kiểm tra.}

\begin{bt}[1,0 điểm]
Tính nhẩm:
\begin{multicols}{2}
\begin{enumerate}[a)]
  \item $9 + 4$
  \item $7 + 5$
  \item $12 - 8$
  \item $14 - 9$
\end{enumerate}
\end{multicols}
\loigiai{
\begin{itemize}
  \item $9 + 4 = 13$
  \item $7 + 5 = 12$
  \item $12 - 8 = 4$
  \item $14 - 9 = 5$
\end{itemize}
}
\end{bt}

\begin{bt}[1,5 điểm]
Đặt tính rồi tính:
\begin{multicols}{3}
\begin{enumerate}[a)]
  \item $23 + 52$
  \item $85 - 26$
  \item $46 + 46$
\end{enumerate}
\end{multicols}
\loigiai{
\begin{itemize}
  \item $23 + 52 = 75$
  \item $85 - 26 = 59$
  \item $46 + 46 = 92$
\end{itemize}
}
\end{bt}

\begin{bt}[1,0 điểm]
Tính giá trị biểu thức kèm đơn vị đo:
\begin{enumerate}[a)]
  \item $31\text{ kg} + 10\text{ kg} - 9\text{ kg}$
  \item $40\text{ }l - 11\text{ }l + 9\text{ }l$
\end{enumerate}
\loigiai{
\begin{enumerate}[a)]
  \item $31\text{ kg} + 10\text{ kg} - 9\text{ kg} = 41\text{ kg} - 9\text{ kg} = 32\text{ kg}$.
  \item $40\text{ }l - 11\text{ }l + 9\text{ }l = 29\text{ }l + 9\text{ }l = 38\text{ }l$.
\end{enumerate}
}
\end{bt}

\begin{bt}[1,5 điểm]
Bạn Giang có $30$ viên bi. Bạn Hà có nhiều hơn bạn Giang $7$ viên bi. Hỏi bạn Hà có bao nhiêu viên bi?
\loigiai{
Số viên bi bạn Hà có là:
\[ 30 + 7 = 37 \text{ (viên bi)} \]
\textbf{Đáp số:} $37$ viên bi.
}
\end{bt}

\begin{bt}[1,0 điểm]
Một thùng chứa đầy $50\text{ }\text{kg}$ đường kính. Buổi sáng cửa hàng đã bán ra $19\text{ }\text{kg}$. Hỏi trong thùng còn lại bao nhiêu ki-lô-gam đường kính?
\loigiai{
Trong thùng còn lại số ki-lô-gam đường kính là:
\[ 50 - 19 = 31\text{ }\text{kg} \]
\textbf{Đáp số:} $31\text{ }\text{kg}$ đường kính.
}
\end{bt}

\begin{bt}[1,0 điểm]
Tìm số tròn chục lớn nhất có hai chữ số mà tổng các chữ số của nó bằng $8$.
\loigiai{
Số tròn chục có chữ số hàng đơn vị là $0$.\\
Chữ số hàng chục là: $8 - 0 = 8$.\\
Vậy số cần tìm là $80$.\\
\textbf{Đáp số:} $80$.
}
\end{bt}
